\section{An infinite prime family for the generalized halving Recam\'an walk}
\label{sec:prime-division-recaman}

OEIS A400939, introduced by Zlatko Damijanic on October 10, 2026,
defines a walk that divides its preceding value by two when that produces
an unvisited integer, and otherwise adds the current index. Its
generalized-divisor conjecture replaces two by an integer that is not a
power of two and predicts eventual permanent addition. We prove an
infinite family of that conjecture, with the stronger conclusion that
exactly one division occurs. The full conjecture and the original
divisor-two occurrence conjecture are not settled here.

For a positive integer $k$, the source rule is
\[
 a_k(1)=1,\qquad
 a_k(n+1)=
 \begin{cases}
  a_k(n)/k,& k\mid a_k(n)\text{ and }
       a_k(n)/k\notin\{a_k(1),\ldots,a_k(n)\},\\
  a_k(n)+n+1,&\text{otherwise},
 \end{cases}
 \qquad n\ge1.
\]
Only the division branch requires an unvisited value. Write
$T_n=n(n+1)/2$.

\paragraph{Theorem.}
Let $p$ be an odd prime with $p\equiv2$ or $3\pmod5$, and put
$q=(p-1)/2$. If $q$ is not a positive triangular number, the unique
division occurs at index $d=p$, and
\[
 a_p(n)=T_n\quad(1\le n<p),\qquad
 a_p(n)=\frac{n(n+1)-p^2-1}{2}\quad(n\ge p).
\]
If $q$ is a positive triangular number, the unique division occurs at
$d=p+1$, and
\[
 a_p(n)=T_n\quad(1\le n\le p),\qquad
 a_p(n)=\frac{n(n+1)-(p+1)^2}{2}\quad(n\ge p+1).
\]
In each case, $p\nmid a_p(n)$ for every $n\ge d$. There are arbitrarily
large primes to which this theorem applies.

\paragraph{Proof.}
If $1\le n$ and $n+1<p$, then $p\nmid T_n$: a contrary divisibility
would imply $p\mid n(n+1)$ although neither positive factor is as large
as $p$. Induction in the literal walk therefore gives
$a_p(n)=T_n$ for $1\le n<p$.

At the next step, the preceding value is
$T_{p-1}=p(p-1)/2=pq$. If $q$ is unvisited, division gives
$a_p(p)=q$, and hence
\[
 2a_p(p)+p^2+1=p(p+1).
\]
If $q$ has appeared, the rule instead gives $a_p(p)=T_p$. The next
proposed quotient is $q+1$. Two positive triangular numbers cannot be
consecutive: for $i>j\ge1$,
\[
 T_i-T_j\ge T_{j+1}-T_j=j+1\ge2.
\]
Thus $q+1$ has not appeared in the triangular prefix, and
$a_p(p+1)=(p+1)/2$. In this case,
\[
 2a_p(p+1)+(p+1)^2=(p+1)(p+2).
\]
This initial-division argument holds for every odd prime.

In the first case let $D=p^2+1$, and in the second let $D=(p+1)^2$.
At the first division index $d$ we have
\[
 2a_p(d)+D=d(d+1),\qquad D\equiv1\pmod p.
\]
Suppose that $2a_p(n)+D=n(n+1)$ for some $n\ge d$. If
$p\mid a_p(n)$, reduction modulo $p$ would give $n(n+1)\equiv1$,
and consequently
\[
 (2n+1)^2=4n(n+1)+1\equiv5\pmod p.
\]
Quadratic reciprocity, since $5\equiv1\pmod4$, identifies the quadratic
character of $5$ modulo $p$ with that of $p$ modulo $5$. The only square
residues modulo $5$ are $0,1,4$, so the residue classes in the theorem
make this impossible. Hence the walk must add $n+1$. The invariant is
preserved because
\[
 2a_p(n+1)+D=n(n+1)+2(n+1)=(n+1)(n+2).
\]
Induction proves the stated tail and its nondivisibility. Together with
the initial dynamics, this proves uniqueness of the division.

Finally, Dirichlet's theorem gives arbitrarily large primes
$p\equiv3\pmod5$. These primes form an infinite subfamily of the
classification. This last step and quadratic reciprocity are credited
classical dependencies.\hfill$\square$

\paragraph{Formalization and source scope.}
The accompanying Lean module \texttt{A400939/Proof.lean} defines the
literal visited-value walk, explicitly excluding its auxiliary zero
index from the visited set. The declaration
\texttt{walk\_followsRule} proves the exact source-rule bridge;
\texttt{first\_prime\_division} checks the first division for all odd
primes; and \texttt{division\_at\_prime\_iff} checks the triangular test
selecting the index. The main theorem
\texttt{walk\_prime\_residue\_classes} proves the full prefix, unique
division, permanent exact tail, and tail nondivisibility.
\texttt{arbitrarily\_large\_single\_division} proves infinitude using
mathlib's existing Dirichlet theorem. Both end theorems depend only on
the standard Lean axioms \texttt{propext}, \texttt{Classical.choice},
and \texttt{Quot.sound}.

The primary statement is pinned to OEIS A400939 revision 58,
November 5, 2026, in OEIS data-export commit
\texttt{df38d77b83a78d8ab5aaa501dec6e3a3db977f01}.
Bounded exact-identifier and semantic searches of both available web
engines, the current trureturing catalog, and its GitHub issues found
no prior proof of this particular prime-family classification. This
documents a bounded check rather than establishing worldwide priority.
The package is an infinite-family partial result, not a claim that
the entire A400939 generalized conjecture has been resolved.
